%! Interpreting p-Values — Unit 6, Lesson 6.5 — Homework
\documentclass[10pt]{article}
\usepackage{apstats-article}
\usepackage{apstats-boxes}
\usepackage{grading}
\gradingassignment{apstats-u06-l05-hw}
\gradingcourse{advanced_statistics}
\newcommand{\TallMath}[1]{$\displaystyle #1\rule[-1.4em]{0pt}{3.2em}$}

\begin{document}

\pageheader{Unit 6, Lesson 6.5}{Homework}

\vspace{0.15in}

\begin{enumerate}

  \item For each set of information, compute the $p$-value.
        \begin{enumerate}[label=\alph*.]
          \item $H_a: p > 0.40$, $z = 1.88$. Tail: \rlap{\smash{\gradedblank[id=hw1a1, type=short, points=1, answer=right, check=text]{2cm}{0.17in}}}\phantom{\blank{2cm}} \quad $p$-value $=$ \rlap{\smash{\gradedblank[id=hw1a2, type=short, points=1, answer=0.0301, check=numeric, tolerance=0.002]{4cm}{0.17in}}}\phantom{\blank{4cm}}
          \item $H_a: p < 0.65$, $z = -2.31$. Tail: \rlap{\smash{\gradedblank[id=hw1b1, type=short, points=1, answer=left, check=text]{2cm}{0.17in}}}\phantom{\blank{2cm}} \quad $p$-value $=$ \rlap{\smash{\gradedblank[id=hw1b2, type=short, points=1, answer=0.0104, check=numeric, tolerance=0.002]{4cm}{0.17in}}}\phantom{\blank{4cm}}
          \item $H_a: p \ne 0.50$, $z = 2.05$. Tail: \rlap{\smash{\gradedblank[id=hw1c1, type=short, points=1, answer=both, check=text]{2cm}{0.17in}}}\phantom{\blank{2cm}} \quad $p$-value $=$ \rlap{\smash{\gradedblank[id=hw1c2, type=short, points=1, answer=0.0404, check=numeric, tolerance=0.002]{4cm}{0.17in}}}\phantom{\blank{4cm}}
        \end{enumerate}

  \item A company claims 80\% of their customers are satisfied. A consumer watchdog group
        believes it's lower. From 250 randomly selected customers, 185 say satisfied.
        \begin{enumerate}[label=\alph*.]
          \item State $H_0$ and $H_a$. \gradedhide{\writeline\\\gradedhere{\gradedblank[id=hw2a, type=rubric, points=2, display={\bfseries $H_0\colon p = 0.80$; $H_a\colon p < 0.80$}]{\linewidth+6pt}{\gradedlinesheight{9}}}\writelines{9}}
          \item $\hat p = $ \rlap{\smash{\gradedblank[id=hw2b1, type=short, points=1, answer=0.74, check=numeric]{2cm}{0.17in}}}\phantom{\blank{2cm}} \quad $z \approx$ \rlap{\smash{\gradedblank[id=hw2b2, type=short, points=2, answer=-2.37, check=numeric, tolerance=0.02]{2cm}{0.17in}}}\phantom{\blank{2cm}} \quad $p$-value $\approx$ \rlap{\smash{\gradedblank[id=hw2b3, type=short, points=2, answer=0.0089, check=numeric, tolerance=0.001]{3cm}{0.17in}}}\phantom{\blank{3cm}}
          \item Write a complete contextual interpretation of the $p$-value. \gradedhide{\writeline\\\gradedhere{\gradedblank[id=hw2c, type=rubric, points=3, display={If 80\% of all customers are satisfied, there is about a 0.89\% chance of a sample proportion of 0.74 or lower (random samples of 250).}]{\linewidth+6pt}{\gradedlinesheight{2}}}\writelines{9}}
        \end{enumerate}

  \newpage
  \item \textbf{Correcting errors.} For each statement, identify and correct the error.
        \begin{enumerate}[label=\alph*.]
          \item ``$p$-value $= 0.03$ means there is only a 3\% chance the study results were
                due to chance.'' \gradedhide{\writeline\\\gradedhere{\gradedblank[id=hw3a, type=rubric, points=2, display={It is not P(results due to chance): IF $H_0$ were true, there is a 3\% chance of a result at least this extreme.}]{\linewidth+6pt}{\gradedlinesheight{1}}}\writelines{9}}
          \item ``Since $p$-value $= 0.62$ is large, we have proven $H_0$ is true.'' \gradedhide{\writeline\\\gradedhere{\gradedblank[id=hw3b, type=rubric, points=2, display={A large $p$-value does not prove $H_0$: the data give no convincing evidence against it, so we fail to reject $H_0$.}]{\linewidth+6pt}{\gradedlinesheight{1}}}\writelines{9}}
        \end{enumerate}

\end{enumerate}

\begin{reflectionbox}
\textbf{Preview:} Lesson 6.6 introduces the significance level $\alpha$ and the formal
reject/fail-to-reject decision rule. You will compare the $p$-value to $\alpha$ to make
a formal conclusion about $H_0$ --- and write that conclusion in context.
\end{reflectionbox}

\end{document}
